% !TeX root = ../../Book3.tex
% 纲要标题：### 11.2 Hilbert 函数
% 纲要位置：工作记录/计划纲要.md，第 2265 行。
% 纲要细目（写作范围）：
% * 分次片的维数或长度
% * Hilbert 级数的有理性
% * 从级数读取增长
% * 用初始单项式理想和包含排除实际计算 Hilbert 数据
% * 按原多项式次数计算同时变号不变环的权二分次 Hilbert 级数，用乘以齐次非零因子的正合列与单项式计数互相核验
% * 证明各次有限维的分次代数之间若有保持次数的满射且 Hilbert 级数相同，则该满射为同构；给出无映射时结论失效的反例
% #### 11.2.1 分次片的维数与长度
% #### 11.2.2 Hilbert 级数的有理性
% #### 11.2.3 Hilbert 级数与增长
% #### 11.2.4 从初始理想计算 Hilbert 数据
% #### 11.2.5 不变环的分次与展示核验

\section{Hilbert 函数}
\label{b3:sec:1102}

对每一个次数分别计数，比只说模“无限维”提供了更多信息。例如 \(k[x]\) 的每个分次片只有一维，而 \(k[x,y]\) 的第 \(n\) 片有 \(n+1\) 维。下面把这些数同时装进一个形式级数；级数的运算只使用系数，不涉及分析中的收敛。

\subsection{分次片的维数与长度}
\label{b3:subsec:1102:01}

本节设 \(A\) 为交换正分次环，\(A_0\) 为 Artin 环，\(A\) 由有限个齐次元素 \(x_1,\ldots,x_r\) 作为 \(A_0\) 代数生成，\(d_i=\deg x_i>0\)，\(M\) 为有限生成分次 \(A\) 模。记号 \(A=A_0[x_1,\ldots,x_r]\) 表示这些元素生成了 \(A\)，允许它们满足关系。给不定元 \(X_i\) 赋次数 \(d_i\)，就有分次满同态
\[
A_0[X_1,\ldots,X_r]\twoheadrightarrow A,
\qquad X_i\longmapsto x_i.
\]
由\cref{b3:prop:graded-noetherian}，\(A\) 是 Noether 环；各 \(M_n\) 是有限生成的 \(A_0\) 模，故长度有限。

\begin{definition}\label{b3:def:hilbert-function-series}
设 \(A=\bigoplus_{n\ge0}A_n\) 为交换正分次环，\(A_0\) 为 Artin 环，\(A=A_0[x_1,\ldots,x_r]\)，其中 \(x_i\) 为正次数齐次代数生成元，不要求代数独立。设 \(M\) 为有限生成分次 \(A\) 模。\(M\) 的\term[hilbert-hanshu]{Hilbert 函数}[Hilbert function]及\term[hilbert-jishu]{Hilbert 级数}[Hilbert series]分别为
\[
h_M(n)=\ell_{A_0}(M_n),\qquad
H_M(t)=\sum_{n\in\mathbb Z}h_M(n)\cdot t^n\in\mathbb Z(\!(t)\!).
\]
这里只有有限多个负次幂。若 \(A_0=k\) 为域，长度就是 \(k\) 维数。
\symidx{\(h_M(n)\)、\(H_M(t)\)：分次 Hilbert 函数与级数}
\end{definition}

短正合列 \(0\rightarrow L\rightarrow M\rightarrow N\rightarrow0\) 若保持次数，则 \(h_M(n)=h_L(n)+h_N(n)\)，从而 \(H_M=H_L+H_N\)。长度可加是因为将子模的合成列与商模合成列的逆像连接，就得到中间项的合成列。次数平移则给出
\begin{equation}\label{b3:eq:hilbert-shift}
H_{M(a)}(t)=\sum_n h_M(n+a)t^n=t^{-a}\cdot H_M(t).
\end{equation}

例如 \(k[x_1,\ldots,x_r]\) 取标准分次时，次数 \(n\) 的单项式数为 \(\binom{n+r-1}{r-1}\)、\(r\ge1\)，所以
\[
H_{k[x_1,\ldots,x_r]}(t)=\dfrac1{(1-t)^r}.
\]
这也可直接把每个变量的几何级数相乘得到。若底环为任意 Artin 环 \(A_0\)，同一自由单项式基使分子乘上 \(\ell_{A_0}(A_0)\)。

\subsection{Hilbert 级数的有理性}
\label{b3:subsec:1102:02}

下面的\term[hilbert-serre-dingli]{Hilbert–Serre 定理}[Hilbert--Serre theorem]把有限生成条件转化为级数分母的有限形式。

\begin{theorem}[Hilbert–Serre]\label{b3:thm:hilbert-serre}
设 \(A=\bigoplus_{n\ge0}A_n\) 为交换正分次环，\(A_0\) 为 Artin 环，\(A=A_0[x_1,\ldots,x_r]\)，其中 \(x_i\) 为次数 \(d_i>0\) 的齐次代数生成元，不要求代数独立。设 \(M\) 为有限生成分次 \(A\) 模。则存在整系数 Laurent 多项式 \(Q(t)\)，使
\[
H_M(t)=\dfrac{Q(t)}{\prod_{i=1}^r(1-t^{d_i})}.
\]
若 \(M_n=0\) 对 \(n<0\) 成立，则可取 \(Q(t)\in\mathbb Z[t]\)。
\end{theorem}
\begin{proof}
对代数生成元个数 \(r\) 归纳。\(r=0\) 时，\(M\) 是有限生成的 \(A_0\) 模，有限组齐次生成元的次数之外分量均为零，所以级数为 Laurent 多项式。

为了对生成元数归纳，要让最后一个生成元 \(x_r\) 的作用变成零。商模 \(L=M/x_rM\) 记录模掉它以后剩下的部分；但乘以 \(x_r\) 未必单射，还须用 \(K=(0:_Mx_r)\) 记录被这次乘法杀掉的部分。令 \(d=d_r\)，乘 \(x_r\) 给出分次正合列
\[
0\longrightarrow K(-d)\longrightarrow M(-d)
\xrightarrow{\ x_r\ }M\longrightarrow L\longrightarrow0.
\]
这里 \(m\in M(-d)_n=M_{n-d}\) 时，\(x_rm\in M_n\)，所以乘法保持次数，源中的核正是 \(K(-d)\)。由 \(A\) 的 Noether 性，\(K,L\) 都有限生成；它们被 \(x_r\) 零化，因而是商代数 \(A/(x_r)\) 上的有限生成分次模。这个商代数由前 \(r-1\) 个生成元的像作为 \(A_0\) 代数生成。逐片长度可加与\eqref{b3:eq:hilbert-shift}给出
\begin{equation}\label{b3:eq:hilbert-induction}
(1-t^d)H_M(t)=H_L(t)-t^d\cdot H_K(t).
\end{equation}
对右边应用归纳假设即得结论。若 \(M\) 没有负次分量，则它的级数及所乘的分母都没有负次幂，因此所得 Laurent 多项式实际上没有负次项。
\end{proof}

这里先证明 \(K\) 有限生成再作归纳，是 Noether 性不可省略的一步。此证明可对照 \cite[\S13, Theorem 13.2, pp. 94--95]{Matsumura1989CommutativeRingTheory}。

\subsection{Hilbert 级数与增长}
\label{b3:subsec:1102:03}

标准分次时，分母全由 \(1-t\) 构成。把公共因子约去，可写
\[
H_M(t)=\dfrac{Q(t)}{(1-t)^d},\qquad Q(1)\ne0
\]
或 \(d=0\)、\(H_M\) 本身为 Laurent 多项式。对于非零模，将约分后分母中 \(1-t\) 的指数 \(d\) 称为在 \(t=1\) 处的极点阶数，没有极点时取 \(d=0\)。这是对应有理函数的代数性质，只须检查公因子的约去，不涉及解析邻域或极限。下一节将证明它恰等于模支集的维数。

\ProseParagraphBreak
例如对 \(A=k[x,y]\)、\(M=A/(x^2)\)，乘 \(x^2\) 在 \(A\) 上单射，故
\[
H_M(t)=\dfrac{1-t^2}{(1-t)^2}=\dfrac{1+t}{1-t}.
\]
所以 \(h_M(0)=1\)，而 \(h_M(n)=2\) 对 \(n\ge1\) 成立。若改取 \(A/(x^2,xy)\)，标准单项式为 \(1,x,y,y^2,\ldots\)，因而
\[
H_{A/(x^2,xy)}(t)=\dfrac1{1-t}+t.
\]
额外的嵌入部分只改变有限多个低次值，却不改变最终的常数一。

\ProseParagraphBreak
在 Hilbert–Serre 定理的假设下，一般正分次已经保证级数有理；标准分次进一步把分母统一为 \(1-t\) 的幂，由此可从其系数得到下一节的最终多项式结论。一般正权分次没有这一保证：上节 \(k[x]\)、\(\deg x=2\) 的级数为 \(1/(1-t^2)\)。若其函数最终由一个多项式表示，无穷多个奇数处的零值会迫使这个多项式恒为零，却与偶数处的值为一矛盾。每个同余类可分别呈现多项式行为，但这不是标准分次结论中的同一个多项式。

\subsection{从初始理想计算 Hilbert 数据}
\label{b3:subsec:1102:04}

Hilbert 级数的有理性还可以变成实际的计数方法。齐次理想的 Gröbner 基把每个分次片的基列成标准单项式，计算由此只涉及哪些单项式被排除。

\begin{proposition}\label{b3:prop:hilbert-initial-count}
设 \(k\) 为域，\(S=k[x_1,\ldots,x_n]\) 取标准分次，\(I\subseteq S\) 为齐次理想，并固定全局单项式序。则 \(S/I\) 与 \(S/\operatorname{in}(I)\) 有相同的 Hilbert 函数。若 \(\operatorname{in}(I)=(m_1,\ldots,m_s)\)，则
\[
H_{S/I}(t)=\dfrac{\displaystyle\sum_{J\subseteq\{1,\ldots,s\}}
(-1)^{|J|}t^{\deg\operatorname{lcm}(m_j:j\in J)}}{(1-t)^n},
\]
空集的最小公倍式约定为一。
\end{proposition}
\begin{proof}
从有限齐次生成组开始作 Buchberger 算法，可得到由齐次多项式组成的 Gröbner 基：S 多项式与每一步非零余式仍齐次。对齐次多项式作除法时，每个乘积保持原次数。因此次数为 \(d\) 的标准单项式的剩余类构成 \((S/I)_d\) 的基；\(I\ne S\) 时，独立性由\cref{b3:prop:standard-basis-finite}的首项论证得到，该论证不要求商空间有限维。\(I=S\) 时，两边的商环均为零，标准单项式集合为空，结论也成立。

计数时，若只有两个单项式生成元 \(m_1,m_2\)，被排除的单项式是“被 \(m_1\) 整除”与“被 \(m_2\) 整除”这两个集合的并。两者的交由最小公倍式的倍数构成，所以被排除部分的级数是
\[
\dfrac{t^{\deg m_1}+t^{\deg m_2}-t^{\deg\operatorname{lcm}(m_1,m_2)}}{(1-t)^n}.
\]
一般地，对每个 \(j\)，被 \(m_j\) 整除的单项式集合的计数级数为 \(t^{\deg m_j}/(1-t)^n\)，若同时被若干生成元整除，则等价于被它们的最小公倍式整除。对每个次数中的有限集合应用包含排除，求出这些集合的并，再从全部单项式中扣除这个并，就得到标准单项式的计数。将各次等式组成形式级数，便得到所示公式，其中空集项贡献全部单项式的级数。
\end{proof}

例如在 \(S=k[x,y,z]\) 中取 \(I=(x^2,xy+y^2)\)，按字典序 \(x\succ y\succ z\) 计算。恒等式
\[
y^3=(y-x)(xy+y^2)+yx^2
\]
表明须补入 \(y^3\)。三对 S 多项式中，第一对依次用 \(xy+y^2\)、\(y^3\) 约化为零，另外两对分别为零或 \(y^4\)，也有零余式。因此 \(x^2,xy+y^2,y^3\) 构成 Gröbner 基，其初始理想为 \((x^2,xy,y^3)\)。标准单项式恰为
\[
z^j,\quad xz^j,\quad yz^j,\quad y^2z^j\qquad(j\ge0).
\]
所以
\[
H_{S/I}(t)=\dfrac{1+2t+t^2}{1-t},\qquad
h_{S/I}(d)=\begin{cases}1,&d=0,\\3,&d=1,\\4,&d\ge2.\end{cases}
\]
下一节的维数与重数定理将说明：分母的一阶极点给出维数一，分子在 \(t=1\) 处的值 \(1+2+1=4\) 给出重数四。这里先通过逐项计数得到函数，随后才解释这些数的几何意义。

\subsection{不变环的分次与展示核验}
\label{b3:subsec:1102:invariant-hilbert}

在\cref{b3:prop:sign-invariant-presentation}中，生成性和关系理想都已由单项式证明。现在保留原多项式的全次数，用 Hilbert 级数重新观察这个环。次数的选择不可略去：\(x^2,xy,y^2\) 的次数都是二，所以对应的三变量多项式环应取权二分次。

\begin{example}\label{b3:exm:sign-invariant-hilbert}
设 \(k\) 为特征不为二的域，\(B=k[x,y]^{C_2}\)，其中非单位元把 \(x,y\) 同时变号。给 \(k[x,y]\) 取 \(\deg x=\deg y=1\)，并令 \(B_n=B\cap k[x,y]_n\)。则
\begin{equation}\label{b3:eq:sign-invariant-hilbert}
H_B(t)=\dfrac{1-t^4}{(1-t^2)^3}
=\dfrac{1+t^2}{(1-t^2)^2},
\qquad
\dim_k B_n=
\begin{cases}
n+1,&n\ge0\;\text{为偶数},\\
0,&n\ge0\;\text{为奇数}.
\end{cases}
\end{equation}
\end{example}
\begin{proof}
该作用保持次数，不变多项式的每个齐次部分仍不变，所以所定义的 \(B_n\) 给出分次。由\cref{b3:prop:sign-invariant-presentation}，令 \(S=k[u,v,w]\)、\(\deg u=\deg v=\deg w=2\)，便有分次同构
\[
B\cong S/(q),\qquad q=uw-v^2,\qquad \deg q=4.
\]
\(S\) 是整环且 \(q\ne0\)，因此乘以 \(q\) 在 \(S\) 上单射，得到保持次数的短正合列
\[
0\longrightarrow S(-4)\xrightarrow{\ q\ }S
\longrightarrow B\longrightarrow0.
\]
每个变量的幂贡献级数 \(1/(1-t^2)\)，故 \(H_S(t)=(1-t^2)^{-3}\)。由逐次维数可加和\eqref{b3:eq:hilbert-shift}，便得
\[
H_B(t)=H_S(t)-t^4\cdot H_S(t)=\dfrac{1-t^4}{(1-t^2)^3}.
\]
这里使用的是 \(q\) 在取商前的 \(S\) 上为非零因子；不能仅从“有一条齐次关系”便省略单射性检查。

另一方面，偶数次 \(2m\) 的全部 \(2m+1\) 个单项式都固定，奇数次没有非零不变元。于是也可直接计算
\[
H_B(t)=\sum_{m\ge0}(2m+1)t^{2m}
=\dfrac{2t^2}{(1-t^2)^2}+\dfrac1{1-t^2}
=\dfrac{1+t^2}{(1-t^2)^2}.
\]
中间等式可逐项乘以 \((1-t^2)^2\) 核验，完全是形式级数恒等式。这同时给出所述 Hilbert 函数。
\end{proof}

由于本例所有奇次分量均为零，另定义 \(\widetilde B_m=B_{2m}\) 不会遗漏任何元素。把原次数二改记为次数一后，三个代数生成元都位于一次分量，得到标准分次环及级数
\[
H_{\widetilde B}(z)=\dfrac{1+z}{(1-z)^2},
\qquad H_B(t)=H_{\widetilde B}(t^2).
\]
底层环、乘法和 Krull 维数不变，Hilbert 函数却按新的次数下标计数。原分次的奇次空片使函数不能最终成为一个多项式，新分次则有 \(h_{\widetilde B}(m)=2m+1\)，可应用标准分次的 Hilbert 多项式与重数结论。直接把 \(u,v,w\) 视为次数一而仍声称计算的是 \(B_n\)，会把所有非零分次片移到错误的位置。若一般分次环有非零奇次分量，只取偶次分量会得到一个子环，不能当作同一底层环的这次重编号。

\ProseParagraphBreak
在上述算例中，第二次计数只是核对已经证明的同构。实际计算时，它也能用来证明尚待核验的展示，但必须先有一个保持次数的满射。

\begin{proposition}\label{b3:prop:hilbert-surjection-isomorphism}
设 \(k\) 为域，\(R=\bigoplus_{n\ge0}R_n\)、\(A=\bigoplus_{n\ge0}A_n\) 为各分次片有限维的分次 \(k\) 代数，\(\psi:R\to A\) 为保持次数的 \(k\) 代数满同态。若 \(H_R(t)=H_A(t)\)，则 \(\psi\) 是同构。
\end{proposition}
\begin{proof}
设 \(K=\ker\psi\)。因为 \(\psi\) 保持次数，\(K\) 是齐次理想。任取 \(a\in A_n\) 并取其原像 \(r=\sum_jr_j\)，比较 \(\psi(r)\) 的各次分量可知 \(\psi(r_n)=a\)，所以每个 \(\psi_n:R_n\to A_n\) 都满射。逐次取核得
\[
0\longrightarrow K_n\longrightarrow R_n
\xrightarrow{\ \psi_n\ }A_n\longrightarrow0.
\]
级数相同意味着 \(\dim_kR_n=\dim_kA_n\)，故 \(K_n=0\) 对每个 \(n\) 成立。因 \(K=\bigoplus_nK_n\)，核为零。
\end{proof}

例如在同时变号作用中，若已证明 \(x^2,xy,y^2\) 生成不变环，并验证 \(uw-v^2\) 是关系，就得到 \(S/(uw-v^2)\to B\) 的分次满射。取商前乘以 \(uw-v^2\) 的正合列计算其源的级数，独立的偶数次单项式计数计算其目标的级数；两者相等后，\cref{b3:prop:hilbert-surjection-isomorphism}便给出核恰为这一条关系的另一证明。Hilbert 级数只记录各次分量的维数，没有记录它们之间的乘法；若尚未证明候选元生成目标，或未说明映射保持所用分次，单独写出两个相同的有理函数还不能得到这个结论。
